% Source: Weinberg GR, physical PDF pp. 51--55
%         (printed pp. 25--29).
% Coverage: Section 2.1, Lorentz Transformations; equations
%           (2.1.1)--(2.1.21).
% Figures/tables/footnotes: no figures or tables; reference markers 1 and 2.
% Status: source-reviewed and compile-clean.
% Uncertainties: none.

\subsection{Lorentz Transformations}\label{sec:2.1}

The Principle of Special Relativity states that the laws of nature are invariant
under a particular group of space-time coordinate transformations, called
Lorentz transformations. We saw at the end of Chapter~1 that Newton's laws of
motion are invariant under the Galilean coordinate transformations
\eqref{eq:1.3.2}, but that Maxwell's equations are not, and that Einstein
resolved this conflict by replacing Galilean invariance with Lorentz invariance.
I shall not continue this discussion in historical terms, but shall simply define
the Lorentz transformations, and then show how Lorentz invariance guides our
search for the laws of nature.

A Lorentz transformation is a transformation from one system of space-time
coordinates \(x^\mu\) to another system \(x'^\mu\), so that
\begin{equation}
  x'^\mu=\Lambda^\mu{}_\nu x^\nu+a^\mu ,
  \tag{2.1.1}\label{eq:2.1.1}
\end{equation}
where \(a^\mu\) and \(\Lambda^\mu{}_\nu\) are constants, restricted by the
conditions
\begin{equation}
  \Lambda^\rho{}_\mu\Lambda^\sigma{}_\nu\eta_{\rho\sigma}
  =\eta_{\mu\nu},
  \tag{2.1.2}\label{eq:2.1.2}
\end{equation}
with
\begin{equation}
  \eta_{\mu\nu}=\diag(-1,+1,+1,+1).
  \tag{2.1.3}\label{eq:2.1.3}
\end{equation}
In our notation \(\mu,\nu,\rho,\ldots\) will always run over the four values
\(0,1,2,3\), with \(x^0=t\) and \(x^1,x^2,x^3\) the Cartesian components of
the position vector \(\mathbf{x}\). We shall use natural units in which the speed
of light is unity, so all \(x^\mu\) have the dimension of length. Any index, like
\(\nu\) in Eq.~\eqref{eq:2.1.1}, that appears twice, once as a subscript and once
as a superscript, is understood to be summed over unless otherwise noted; that
is, Eq.~\eqref{eq:2.1.1} is an abbreviation for
\[
  x'^\mu=\Lambda^\mu{}_0x^0+\Lambda^\mu{}_1x^1
  +\Lambda^\mu{}_2x^2+\Lambda^\mu{}_3x^3+a^\mu .
\]

The fundamental property that distinguishes the Lorentz transformations is that
they leave invariant the ``proper time'' \(\dd\tau\), defined by
\begin{equation}
  \dd\tau^2=\dd t^2-\dd\mathbf{x}^2
  =-\eta_{\mu\nu}\,\dd x^\mu \dd x^\nu .
  \tag{2.1.4}\label{eq:2.1.4}
\end{equation}
In a new coordinate system \(x'^\mu\), the coordinate differentials are given
by Eq.~\eqref{eq:2.1.1} as
\[
  \dd x'^\mu=\Lambda^\mu{}_\nu\,\dd x^\nu ,
\]
so the new proper time will be
\[
\begin{aligned}
  \dd\tau'^2
  &=-\eta_{\mu\nu}\,\dd x'^\mu \dd x'^\nu \\
  &=-\eta_{\mu\nu}\Lambda^\mu{}_\rho\Lambda^\nu{}_\sigma
       \dd x^\rho \dd x^\sigma \\
  &=-\eta_{\rho\sigma}\,\dd x^\rho \dd x^\sigma ,
\end{aligned}
\]
and therefore
\begin{equation}
  \dd\tau'^2=\dd\tau^2 .
  \tag{2.1.5}\label{eq:2.1.5}
\end{equation}
It is this property that accounts for the observation by Michelson and Morley
that the speed of light is the same in all inertial systems. A light wave front
will have \(\lvert \dd\mathbf{x}/\dd t\rvert\) equal to the speed of light, which in
our units is unity; hence the propagation of light is described by the statement
that
\begin{equation}
  \dd\tau=0 .
  \tag{2.1.6}\label{eq:2.1.6}
\end{equation}
Performing a Lorentz transformation does not change \(\dd\tau\), so
\(\dd\tau'^2=0\), and therefore
\(\lvert \dd\mathbf{x}'/\dd t'\rvert=1\); that is, the speed of light in the new
coordinate system is still unity.

We can also show that the Lorentz transformations \eqref{eq:2.1.1} are the
\emph{only} nonsingular coordinate transformations \(x\to x'\) that leave
\(\dd\tau^2\) invariant. (Nonsingular means that \(x'(x)\) and \(x(x')\) are
well-behaved differentiable functions, so that the matrix
\(\partial x'^\mu/\partial x^\nu\) has a well-defined inverse
\(\partial x^\nu/\partial x'^\mu\).) A general coordinate transformation
\(x\to x'\) will change \(\dd\tau\) into \(\dd\tau'\), given by
\[
\begin{aligned}
  \dd\tau'^2
  &=-\eta_{\mu\nu}\,\dd x'^\mu \dd x'^\nu \\
  &=-\eta_{\mu\nu}
    \frac{\partial x'^\mu}{\partial x^\rho}
    \frac{\partial x'^\nu}{\partial x^\sigma}
    \dd x^\rho \dd x^\sigma .
\end{aligned}
\]
If this is equal to \(\dd\tau^2\) for all \(\dd x^\rho\), we must have
\begin{equation}
  \eta_{\rho\sigma}
  =\eta_{\mu\nu}
    \frac{\partial x'^\mu}{\partial x^\rho}
    \frac{\partial x'^\nu}{\partial x^\sigma}.
  \tag{2.1.7}\label{eq:2.1.7}
\end{equation}
Differentiation with respect to \(x^\lambda\) gives
\[
  0=\eta_{\mu\nu}
  \frac{\partial^2x'^\mu}{\partial x^\rho\partial x^\lambda}
  \frac{\partial x'^\nu}{\partial x^\sigma}
  +\eta_{\mu\nu}
  \frac{\partial x'^\mu}{\partial x^\rho}
  \frac{\partial^2x'^\nu}{\partial x^\sigma\partial x^\lambda}.
\]
To solve for the second derivatives, we add to this the same equation with
\(\rho\) and \(\lambda\) interchanged, and subtract the same with \(\lambda\)
and \(\sigma\) interchanged; that is,
\[
\begin{aligned}
0=\eta_{\mu\nu}\Bigg[
&\frac{\partial^2x'^\mu}{\partial x^\rho\partial x^\lambda}
 \frac{\partial x'^\nu}{\partial x^\sigma}
+\frac{\partial^2x'^\nu}{\partial x^\sigma\partial x^\lambda}
 \frac{\partial x'^\mu}{\partial x^\rho}\\
&+\frac{\partial^2x'^\mu}{\partial x^\lambda\partial x^\rho}
 \frac{\partial x'^\nu}{\partial x^\sigma}
+\frac{\partial^2x'^\nu}{\partial x^\sigma\partial x^\rho}
 \frac{\partial x'^\mu}{\partial x^\lambda}\\
&-\frac{\partial^2x'^\mu}{\partial x^\rho\partial x^\sigma}
 \frac{\partial x'^\nu}{\partial x^\lambda}
-\frac{\partial^2x'^\nu}{\partial x^\lambda\partial x^\sigma}
 \frac{\partial x'^\mu}{\partial x^\rho}
\Bigg].
\end{aligned}
\]
The last term cancels the second, the penultimate cancels the fourth (because
\(\eta_{\mu\nu}=\eta_{\nu\mu}\)), and the first equals the third, so we are
left with
\[
  0=2\eta_{\mu\nu}
  \frac{\partial^2x'^\mu}{\partial x^\rho\partial x^\lambda}
  \frac{\partial x'^\nu}{\partial x^\sigma}.
\]
But both \(\eta_{\mu\nu}\) and
\(\partial x'^\nu/\partial x^\sigma\) are nonsingular matrices, so this
immediately yields
\begin{equation}
  0=\frac{\partial^2x'^\mu}
  {\partial x^\rho\partial x^\lambda}.
  \tag{2.1.8}\label{eq:2.1.8}
\end{equation}
The general solution of Eq.~\eqref{eq:2.1.8} is of course just the linear
function \eqref{eq:2.1.1}, and by inserting Eq.~\eqref{eq:2.1.1} in
Eq.~\eqref{eq:2.1.7} we see that \(\Lambda^\mu{}_\nu\) must be subject to the
condition \eqref{eq:2.1.2}. This proof is an elementary example of the sort of
thing we do in Chapter~13, on symmetric spaces. (Incidentally, if we had only
assumed that the transformations \(x\to x'\) leave \(\dd\tau\) invariant when
\(\dd\tau=0\), that is, for a particle moving at the speed of light, then we would
have found that these transformations are in general nonlinear, and form a
15-parameter group, the conformal group, which contains the Lorentz
transformations as a subgroup. But the statement that a free particle moves at
constant velocity would not be an invariant statement unless the velocity were
that of light, and since there are massive particles in the world, we must reject
the conformal group as a possible invariance of nature.)

The set of all Lorentz transformations of the form \eqref{eq:2.1.1} is
correctly called the \emph{inhomogeneous Lorentz group}, or the
\emph{Poincar\'e group}. The subset with \(a^\mu=0\) is called the
\emph{homogeneous Lorentz group}. Both the homogeneous and the inhomogeneous
Lorentz groups have subgroups called the \emph{proper} homogeneous and
inhomogeneous Lorentz groups, defined by imposing on
\(\Lambda^\mu{}_\nu\) the additional requirements
\begin{equation}
  \Lambda^0{}_0\geq 1,\qquad \det\Lambda=+1 .
  \tag{2.1.9}\label{eq:2.1.9}
\end{equation}
Note from Eq.~\eqref{eq:2.1.2} that
\begin{equation}
  \bigl(\Lambda^0{}_0\bigr)^2
  =1+\sum_{i=1}^{3}\bigl(\Lambda^i{}_0\bigr)^2\geq1
  \tag{2.1.10}\label{eq:2.1.10}
\end{equation}
and
\begin{equation}
  (\det\Lambda)^2=1 .
  \tag{2.1.11}\label{eq:2.1.11}
\end{equation}
[Equation \eqref{eq:2.1.10} follows upon setting \(\mu=\nu=0\) in
Eq.~\eqref{eq:2.1.2}. Equation \eqref{eq:2.1.11} is derived by writing
Eq.~\eqref{eq:2.1.2} as a matrix equation
\(\eta=\Lambda^{\mathrm T}\eta\Lambda\) and taking its determinant.]
It follows that any \(\Lambda^\mu{}_\nu\) that can be converted to the identity
\(\delta^\mu{}_\nu\) by a continuous variation of its parameters must be a
proper Lorentz transformation, because it is impossible by a continuous change
of parameters to jump from \(\Lambda^0{}_0\leq-1\) to
\(\Lambda^0{}_0\geq+1\), or from \(\det\Lambda=-1\) to
\(\det\Lambda=+1\), and the identity has \(\Lambda^0{}_0=+1\) and
\(\det\Lambda=+1\). The improper Lorentz transformations involve either space
inversion (\(\det\Lambda=-1,\ \Lambda^0{}_0\geq1\)), which is now known not
to be an exact symmetry of nature,\textsuperscript{\hyperref[ch2-ref-1]{1}} or
time reversal (\(\det\Lambda=-1,\ \Lambda^0{}_0\leq-1\)), which is strongly
suspected to be not an exact symmetry of
nature,\textsuperscript{\hyperref[ch2-ref-2]{2}} or their product. We are
dealing almost exclusively with proper Lorentz transformations, and unless
otherwise noted any Lorentz transformation is assumed to satisfy
Eq.~\eqref{eq:2.1.9}.

The proper homogeneous Lorentz transformations have a further subgroup,
consisting of the rotations, for which
\[
  \Lambda^i{}_j=R_{ij},\qquad
  \Lambda^i{}_0=\Lambda^0{}_j=0,\qquad
  \Lambda^0{}_0=1 ,
\]
where \(R_{ij}\) is a unimodular orthogonal matrix (i.e.,
\(\det R=1\) and \(R^{\mathrm T}R=1\)) and the indices \(i,j\) run over the
values \(1,2,3\). With regard to both rotations and the space-time translations
\(x^\mu\to x^\mu+a^\mu\), there is no difference between the Lorentz group and
the Galileo group discussed in Chapter~1. The difference arises only in those
transformations, called \emph{boosts}, that change the velocity of the coordinate
frame.

Suppose that one observer \(O\) sees a particle at rest, and a second observer
\(O'\) sees it moving with velocity \(\mathbf v\). From Eq.~\eqref{eq:2.1.1}
we have
\begin{equation}
  \dd x'^\mu=\Lambda^\mu{}_\nu\,\dd x^\nu ,
  \tag{2.1.12}\label{eq:2.1.12}
\end{equation}
or, since \(\dd\mathbf{x}\) vanishes,
\begin{equation}
  \dd x'^i=\Lambda^i{}_0\,\dd t\qquad(i=1,2,3)
  \tag{2.1.13}\label{eq:2.1.13}
\end{equation}
and
\begin{equation}
  \dd t'=\Lambda^0{}_0\,\dd t .
  \tag{2.1.14}\label{eq:2.1.14}
\end{equation}
Dividing \(\dd\mathbf{x}'\) by \(\dd t'\) gives the velocity \(\mathbf v\), so
\begin{equation}
  \Lambda^i{}_0=v_i\Lambda^0{}_0 .
  \tag{2.1.15}\label{eq:2.1.15}
\end{equation}
We can get a second relation between \(\Lambda^i{}_0\) and
\(\Lambda^0{}_0\) by setting \(\mu=\nu=0\) in Eq.~\eqref{eq:2.1.2}:
\begin{equation}
  -1=\Lambda^\rho{}_0\Lambda^\sigma{}_0\eta_{\rho\sigma}
  =\sum_{i=1}^{3}\bigl(\Lambda^i{}_0\bigr)^2
   -\bigl(\Lambda^0{}_0\bigr)^2 .
  \tag{2.1.16}\label{eq:2.1.16}
\end{equation}
The solution of Eqs.~\eqref{eq:2.1.15} and \eqref{eq:2.1.16} is
\begin{equation}
  \Lambda^0{}_0=\gamma ,
  \tag{2.1.17}\label{eq:2.1.17}
\end{equation}
\begin{equation}
  \Lambda^i{}_0=\gamma v_i ,
  \tag{2.1.18}\label{eq:2.1.18}
\end{equation}
where
\begin{equation}
  \gamma\equiv(1-\mathbf v^2)^{-1/2}.
  \tag{2.1.19}\label{eq:2.1.19}
\end{equation}
The other \(\Lambda^\mu{}_\nu\) are not uniquely determined, because if
\(\Lambda^\mu{}_\nu\) carries a particle from rest to velocity \(\mathbf v\),
then so does \(\Lambda^\mu{}_\rho R^\rho{}_\nu\), where \(R\) is an arbitrary
rotation. One convenient choice that satisfies Eq.~\eqref{eq:2.1.2} is
\begin{equation}
  \Lambda^i{}_j=\delta_{ij}
  +v_iv_j\frac{\gamma-1}{\mathbf v^2},
  \tag{2.1.20}\label{eq:2.1.20}
\end{equation}
\begin{equation}
  \Lambda^0{}_j=\gamma v_j .
  \tag{2.1.21}\label{eq:2.1.21}
\end{equation}
It can easily be seen that any proper homogeneous Lorentz transformation may be
expressed as the product of a boost \(\Lambda(\mathbf v)\) times a rotation
\(R\).
